\hwpchapter{5}{창의성 및 차별성}{10}

\hwpsection{5.1 생산운영 특성을 반영한 예측 성능 개선}

C-BAT는 시간별 생산계획에서 생산 여부와 정규화 로그 생산량을 구분해 입력하고, 가동유형별 기본 곡선, 최근 기준일 곡선, 생산 전달항에 가동·비가동 상태별 날짜 효과와 AR(1) 슬롯 잡음을 더해 하루 경로의 분포를 만든다(2.1절). 분위수, 일 피크 분포, 초과 확률이 모두 같은 경로에서 나온다. 잡음을 가동 상태별로 나눈 것은 1.2절에서 본 비가동일과 가동일의 부하 차이를 반영한 설계이다.

개발 구간과 9월 시험 구간을 합친 합산 평가(피크 평가일 65일)에서 일 피크 MAE는 C-BAT 8.93 kW, 기존 BAT 13.09 kW, M2 14.54 kW였다. 기존 BAT보다 31.8\%, M2보다 38.6\% 낮고, 날짜 단위 짝지은 부트스트랩의 95\% 구간이 두 차이 모두 0을 포함하지 않는다(표~\ref{tab:ch5-pooled}). 개발 구간에서는 비가동일의 일 피크 MAE가 C-BAT 2.50 kW, 기존 BAT 24.11 kW였고, 고피크일 13일의 일 피크 MAE는 C-BAT 4.37 kW, 기존 BAT 12.47 kW, LightGBM(튜닝) 18.27 kW였다. 기존 BAT는 비가동일의 피크를 높게, LightGBM(튜닝)은 고피크일의 피크를 낮게 보았고 C-BAT는 두 실패를 모두 줄였다. 고피크일 차이의 95\% 구간은 기존 BAT 대비 $-8.10$ kW $[-11.19,\ -5.80]$, LightGBM(튜닝) 대비 $-13.90$ kW $[-19.01,\ -9.86]$이다.

\begin{table}[htbp]
\centering
\hwptablefont
\setlength{\tabcolsep}{4pt}
\caption{합산 평가의 C-BAT와 비교 모델(피크 평가일 65일, 6,358슬롯, kW). 차이는 C-BAT $-$ 비교 모델이며 대괄호는 날짜 단위 짝지은 부트스트랩 95\% 구간이다.}\label{tab:ch5-pooled}
\begin{tabular}{llrrl}
\toprule
지표 & 비교 모델 & C-BAT & 비교 모델 & 차이 \\
\midrule
일 피크 MAE & 기존 BAT & 8.93 & 13.09 & $-4.17$ [$-6.95$, $-1.31$] \\
 & M2 & 8.93 & 14.54 & $-5.61$ [$-10.24$, $-1.38$] \\
 & B0 & 8.93 & 34.63 & $-25.70$ [$-42.10$, $-11.43$] \\
\midrule
슬롯 MAE & 기존 BAT & 8.22 & 7.61 & $+0.61$ [$+0.07$, $+1.18$] \\
 & M2 & 8.22 & 10.57 & $-2.35$ [$-4.47$, $-0.42$] \\
 & B0 & 8.22 & 22.84 & $-14.61$ [$-24.17$, $-6.32$] \\
\bottomrule
\end{tabular}
\end{table}
\hwpnote{합산 비교는 C-BAT, 기존 BAT, M2, B0 네 모델에 한정한다. B0\_kind와 LightGBM(튜닝)은 9월 결과가 없어 합산 평가에 넣지 않았다.}

개선되지 않은 부분도 함께 보고한다.
\begin{hwplist}
\item 슬롯 MAE: 합산 평가에서 기존 BAT(7.61 kW)가 C-BAT(8.22 kW)보다 0.61 kW 낮다. 차이의 95\% 구간이 0을 포함하지 않는다.
\item 전체 일 피크: 개발 구간에서 B0\_kind(6.67 kW)가 C-BAT(8.93 kW)보다 낮다(차이 $+2.26$ kW, 95\% 구간 $[0.36,\ 4.12]$). 고피크일에서는 C-BAT(4.37 kW)가 B0\_kind(7.54 kW)보다 낮지만 차이 $-3.17$ kW의 날짜 단위 구간 $[-7.04,\ 0.40]$은 0을 포함한다.
\item 9월 단독: M2의 슬롯 MAE 6.60 kW, 일 피크 MAE 5.53 kW가 C-BAT(7.18 kW, 8.93 kW)보다 낮다. 9월에는 고피크일이 없어 이 조건의 성능은 새로 시험하지 못했다.
\end{hwplist}

구성요소의 기여는 절제 비교로 일부만 확인하였다. 학습형 어텐션 모형에서 잡음을 상태별에서 공통으로 바꾸면 개발 구간의 일 피크 MAE가 9.16 kW에서 15.29 kW로, 고피크일 일 피크 MAE가 4.38 kW에서 7.80 kW로 커졌다. 그러나 두 모형의 최대 split-$\hat R$가 1.308과 2.303으로 수렴 기준 1.05를 넘어, 이 차이를 잡음 설계만의 효과 크기나 기존 BAT 대비 고피크일 개선(8.10 kW)의 기여도로 나누지 않았다. 고정 어텐션은 학습하는 모수(상수 제외)를 31개에서 22개로 줄이며, 학습형보다 슬롯 MAE가 0.17 kW 나쁘고 일 피크 MAE는 0.23 kW 좋다. 기존 BAT에서 생산 전달항을 빼도 슬롯 MAE 7.73 kW, 일 피크 MAE 14.53 kW로 기존 BAT(7.82 kW, 14.50 kW)와 거의 같아 전달항 단독의 기여는 확인하지 못했다. 따라서 개선은 구성 전체의 결과로 읽고 구성요소별로 분해하지 않는다.

\hwpsection{5.2 오류 유형에 따른 보완 대상과 평가 절차}

3장의 오류분석에서 C-BAT의 보완 대상을 오류 유형별로 구분하였다.
\begin{hwplist}
\item 평범한 가동일의 일 피크 과대예측: 개발 구간 가동일 37일의 일 피크 편향은 $+9.74$ kW, 일 피크 MAE는 11.36 kW로 비가동일(2.50 kW)과 고피크일(4.37 kW)보다 크다. 월 최대 갱신 경보의 오경보 9건 가운데 6건(7월 7--16일)도 같은 경향과 겹쳤다(4.2절).
\item 생산이 없는 슬롯의 전력 수준: 가동유형 2(13--19시간)에서 계획상 생산이 없는 슬롯의 슬롯 MAE는 20.7 kW, 편향은 $+13.7$ kW였다. 생산 시작 전이나 종료 뒤의 전력을 가동 수준으로 높게 본다.
\item 7월 종일 가동일의 고피크 미탐지: C90 경보에서 미탐지(seed 평균 4.33일)는 모두 7월의 가동유형 3(20시간 이상) 날이었다. 여름철 부하 상승 폭을 다 따르지 못한 것으로 본다. C90 오경보(seed 평균 8일)는 평범한 가동일의 과대예측에서 나왔다.
\end{hwplist}
세 유형은 원인이 달라 한 가지 보정으로 함께 줄이기 어렵다. 순위와 별개로 취약 조건을 구분했으므로 추가 점검과 경보 보정의 대상을 구체화할 수 있다.

평가 절차는 다음 세 단계로 설계하였다.
\begin{hwplist}
\item 누수 진단: 복사일 122일(완전 복사일 115일, 편집 복사일 7일)이 섞인 자료에서 날짜를 무작위로 나누면 시험일의 사본이 학습에 들어간다. 9월 이전 241일에서 M2의 MAE는 무작위 5-fold 14.91 kW, LOEO 16.64 kW로 차이가 $-1.73$ kW(95\% 구간 $[-3.63,\ -0.04]$)였고, 기존 BAT는 $+0.05$ kW($[-0.22,\ +0.42]$)였다. 이 결과로 복사일을 목표에서 뺀 시간순 검증을 택하였다. 이 값은 누수의 정도를 보는 진단이며 전방 예측 성능의 추정치가 아니다.
\item 사전 등록: 기준일 규칙, 잡음 변형, 위험 확률의 발행 방식, 월 최대 갱신 경보 규칙, 반사실 검증의 채택 기준을 결과를 보기 전에 문서로 고정한 뒤 평가하였다.
\item 확장 실험: C-BAT를 확정한 뒤 사전 등록한 확장 실험 11가지(잡음 변형 4, 기준일 변경 4, 전용 피크 계층 2, M2와의 점예측 앙상블 1)를 개발 구간에서만 평가하였고 채택된 것은 없다. 예를 들어 기준일을 바꾼 네 실험은 슬롯 MAE를 7.63--7.92 kW로 낮췄으나 일 피크 MAE가 9.06--11.93 kW로 커졌다. 평범한 가동일의 과대예측을 줄이면서 고피크일 성능을 유지하는 개선안은 찾지 못하였다. 이는 시험한 입력과 모형의 한계이며, 생산계획에 두 부류를 구분할 정보가 없다는 증거는 아니다.
\end{hwplist}
\hwpnote{9월 시험 구간은 C-BAT를 동결한 뒤 한 번 예측하였고 모델 선택에 쓰지 않았다. 다만 이전 개발 단계에서 한 번 열린 적이 있어 완전히 독립된 홀드아웃은 아니다.}

\hwpsection{5.3 요금 구조에 맞춘 현장 활용}

기본요금은 매달 요금 시간대 최대수요와 래칫 바닥 가운데 큰 값으로 정해져, 요금을 움직이는 날은 월 최대를 새로 쓰는 소수의 날이다(4.2절). 일 피크의 평균 오차가 작다는 것만으로는 요금 관점의 쓸모를 말할 수 없고, 같은 예측에서 갱신 위험이 큰 날을 따로 가려내는 기능이 필요하다. C-BAT는 일 피크 분포에서 일 피크가 C50·C75·C90을 넘을 확률(피크 기준 초과 경보, 3장)을 내고, 요금 시간대 슬롯만 모아 월 최대 갱신 확률과 기대 초과 기본요금(월 최대 갱신 경보, 4.2절)을 낸다. 두 경보는 기준이 다르다.

C-BAT의 차별점은 새로운 알고리즘이 아니라 조합에 있다. 시간별 생산계획을 입력으로 받는 해석 가능한 구조에서 하루 경로의 분포를 만들고, 그 분포를 기본요금이 정해지는 방식에 맞추어 쓴다. 표~\ref{tab:ch5-compare}는 이 과제에서 쓸 수 있는 기존 접근과 C-BAT를 다섯 가지 기능으로 비교한다.

\begin{table}[htbp]
\centering
\hwptablefont
\setlength{\tabcolsep}{4pt}
\caption{기존 접근과 C-BAT의 기능 비교(○ 제공, △ 부분 제공 또는 미적용, $\times$ 없음).}\label{tab:ch5-compare}
\begin{tabular}{L{3.8cm}ccccc}
\toprule
접근 & 시간별 계획 & 예측분포·경로 & 피크 위험확률 & 월 최대 경보 & 해석 가능성 \\
\midrule
유사일·B0 기준일 & $\times$ & $\times$ & $\times$ & $\times$ & ○ \\
GAM·평활 가법모형 & △ & △ & △ & $\times$ & ○ \\
LightGBM(M2) & ○ & △ & △ & $\times$ & △ \\
기존 BAT & ○ & ○ & ○ & △ & ○ \\
C-BAT & ○ & ○ & ○ & ○ & ○ \\
\bottomrule
\end{tabular}
\end{table}
\hwpnote{표시는 각 접근의 일반적인 형태에 대한 판단이며 기능의 유무를 비교한 것이지 성능의 우열이 아니다. GAM은 이 보고서에서 실행하지 않았다. 기존 BAT의 월 최대 경보 △는 경로를 만들 수 있으나 이 보고서에서 적용하지 않았다는 뜻이다.}

경로를 하나로 두면 슬롯 분위수, 일 피크 분포, 피크 기준 초과 확률, 요금 시간대 최댓값의 분포와 초과 기본요금을 같은 표본에서 계산할 수 있다. 이 경보는 개발 구간 중 요금 시간대 관측이 있는 44일에서 유일한 월 최대 갱신일인 7월 19일을 잡았고 B0\_kind는 이날을 놓쳤다(4.2절). 사건이 1건이어서 성능 추정이 아니라 운용 사례이다. 133,120원은 완벽한 조치를 가정한 금액 상한이며, 주 규칙의 가정대로 점검 10회에 회당 3만 원을 쓰면 30만 원으로 상한보다 커서 당월 순편익이 양수라는 근거는 없다(4.3절). 모형의 반사실 일정 이동 절감액은 관측 유사일 비교에서 지지되지 않아 주장하지 않았다(4.4절). 이 보고서는 모형 출력을 관측으로 확인한 범위에서만 현장 제안에 썼다.

\hwpnote{한계: 평가 범위는 한 공장의 2021년 자료이며 다른 계절과 다른 공장으로의 일반화는 확인하지 않았다. 생산실적을 사전 생산계획으로 간주하였고, 합산 비교는 4개 모델에 한정하며, 9월은 이전에 한 번 열렸고 고피크일이 없었다. 개발 구간은 여러 모델 비교에 반복해서 쓰여 선택과 조건별 비교는 탐색적이다. 월 최대 갱신 사례는 1건이고 조치의 실제 효과는 현장 검증이 필요하다.}
